\section{稳定曲线的定义}
\label{sec:definition-stable-curves}
\stfhIndex{I133}{shen^gu^xiao^ye}{深谷--小野（Fukaya–Ono）}
\stfhIndex{I179}{kong^cai^wei^qi}{孔采维奇（Kontsevich）}
\stfhIndex{I256}{wen^ding^qu^xian}{稳定曲线（stable curve）}
\stfhIndex{I257}{wen^ding^ying^she}{稳定映射（stable map）}

本节紧随深谷贤治与小野薫 (FOn99) 的论述。特别地，我们将解释其中定义的\emph{稳定映射拓扑}，并紧随他们的方案给出\emph{紧性与豪斯多夫性}的详细证明，同时作出一些进一步澄清。

\subsection{稳定曲线的光滑描述}
\label{subsec:smooth-description-stable-curves}

先引入闭曲线之稳定映射的定义域。考虑由给定有限指标集 $V$ 标号的紧黎曼曲面 $(\Sigma_v,j_v)$，指标集的元素记为 $v$。形成有限不交并
$\coprod_{v\in V}\Sigma_v$。接着考虑由无序偶对
$e=\{x_e,y_e\}$ 构成的有限集
\[
 E=\left\{\{x_e,y_e\}\mathrel{\Big|}
 x_e\ne y_e\in\coprod_{v\in V}\Sigma_v\right\},
\]
满足
\begin{equation}
 \{x_e,y_e\}\cap\{x_{e'},y_{e'}\}=\varnothing
 \qquad\text{若 }e\ne e'.
 \label{eq:nodal-pairwise-disjoint}
\end{equation}
在 $\coprod_{v\in V}\Sigma_v$ 上定义等价关系 $\sim$：对任意
$x\in\coprod_{v\in V}\Sigma_v$ 令 $x\sim x$，并且当
$\{x_e,y_e\}\in E$ 时令 $x_e\sim y_e$；取带商拓扑的商空间
\[
 \Sigma=\coprod_{v\in V}\Sigma_v/\sim.
\]
称点 $[x_e]=[y_e]$ 为 $\Sigma$ 的\emph{二重点}，并以
\stfhIndex{I105}{er^chong^dian}{二重点（double point）}
$\operatorname{Sing}(\Sigma)$ 表示二重点集合。

为了把 $\Sigma$ 作为（奇异）复簇而赋予复结构，对每个二重点
$[x_e]=[y_e]$ 的一个邻域，给出到 $\mathbb C^2$ 中原点附近局部图形 $xy=0$ 的同胚，并把 $[x_e]$ 映到原点。然后拉回由
$\mathbb C^2$ 诱导在图形
\[
 \{(x,y)\in\mathbb C^2\mid xy=0\}
\]
上的标准复结构。把 $\Sigma$ 上这个复结构记为 $j=j_\Sigma$，并称其为诸 $j_v$ 的\emph{粘合复结构}。由此得到一条奇异复代数曲线。

\begin{definition}\label{def:nodal-riemann-surface}
\stfhIndex{I213}{jie^dian^li^man^qu^mian}{节点黎曼曲面（nodal Riemann surface）}
按上述方式定义的偶对 $(\Sigma,j)$ 称为\emph{节点黎曼曲面}。称每个分支
$(\Sigma_v,j_v)\hookrightarrow(\Sigma,j)$ 为一个\emph{不可约分支}。
\end{definition}

从现在起，在不会混淆时，只用 $\Sigma$ 表示偶对 $(\Sigma,j)$。依构造，$\Sigma$ 具有下列性质：
\begin{enumerate}[label=(\alph*)]
 \item $\pi_v:\Sigma_v\to\Sigma$ 是连续映射（正规化）；
 \item 在每个二重点附近，$\Sigma$ 局部形如 $\mathbb C^2$ 中方程
 $xy=0$ 的图形；
 \item $\Sigma$ 只允许具有 (b) 中给出的奇点类型。
\end{enumerate}
注意，性质 (c) 来自条件\eqref{eq:nodal-pairwise-disjoint}与性质 (b)。

\begin{definition}\label{def:pre-stable-curve}
\stfhIndex{I201}{dai^biao^ji^yu^wen^ding^qu^xian}{带标记预稳定曲线（marked pre-stable curve）}
\stfhIndex{I229}{yu^wen^ding^qu^xian}{预稳定曲线（pre-stable curve）}
\emph{预稳定曲线}是偶对 $(\Sigma,\mathbf z)$，其中
$\Sigma=(\Sigma,j)$ 是节点黎曼曲面，而
$\mathbf z=\{z_1,\ldots,z_m\}\subset\Sigma$ 是满足
\begin{enumerate}[label=(\alph*),start=4]
 \item 当 $i\ne j$ 时 $z_i\ne z_j$；
 \item $z_i\in\Sigma\setminus\operatorname{Sing}(\Sigma)$
\end{enumerate}
的有限集。称 $\mathbf z=\{z_1,\ldots,z_m\}$ 中的点为\emph{标记点}。
\end{definition}

接着在预稳定曲线集合上引入等价关系。

\begin{definition}\label{def:isomorphism-pre-stable-curves}
\stfhIndex{I171}{tong^gou}{同构（isomorphism）}
设 $(\Sigma,\mathbf z)$ 与 $(\Sigma',\mathbf z')$ 是两条（带标记的）预稳定曲线。若连续映射 $\phi:\Sigma\to\Sigma'$ 是同胚，并且 $\phi$ 在每个 $\Sigma_v$ 上的限制可提升为到 $\Sigma'$ 某个不可约分支 $\Sigma_w'$ 上的双全纯映射，则称 $\phi$ 为同构。若还满足 $\phi(z_i)=z_i'$，则称
$\phi:(\Sigma,\mathbf z)\to(\Sigma',\mathbf z')$ 为预稳定曲线的同构。
\end{definition}

自同构 $\phi:(\Sigma,\mathbf z)\to(\Sigma,\mathbf z)$ 称为
\stfhIndex{I057}{zi^tong^gou}{自同构（automorphism）}
$(\Sigma,\mathbf z)$ 的自同构。以 $\operatorname{Aut}(\Sigma,\mathbf z)$ 表示其自同构集合。

\begin{definition}\label{def:stable-curve-finite-automorphisms}
\stfhIndex{I255}{wen^ding}{稳定（stable）}
若预稳定曲线 $(\Sigma,\mathbf z)$ 的
$\#\operatorname{Aut}(\Sigma,\mathbf z)$ 有限，则称它\emph{稳定}；否则称它\emph{不稳定}。
\end{definition}

从几何上看，$(\Sigma,\mathbf z)$ 的稳定性等价于
$\Sigma\setminus\{z_1,\ldots,z_m\}$ 上存在双曲度量。

下面的例子给出稳定性判据的几何描述。回忆一致化定理蕴含
$\Sigma=\mathbb{CP}^1$ 具有唯一的共形结构（或复结构）。

\begin{example}\label{ex:stability-automorphism-dimensions}
\begin{enumerate}[label=(\arabic*)]
 \item 对 $\Sigma=\mathbb{CP}^1$，由于
 $\operatorname{Aut}(\mathbb{CP}^1)=PSL(2,\mathbb C)$，有
 $\dim_{\mathbb C}\operatorname{Aut}(\mathbb{CP}^1)=3$。此外，
 \[
 \dim_{\mathbb C}\operatorname{Aut}(\Sigma,z_1)=2,
 \quad
 \dim_{\mathbb C}\operatorname{Aut}(\Sigma,z_1,z_2)=1,
 \quad
 \dim_{\mathbb C}\operatorname{Aut}(\Sigma,z_1,z_2,z_3)=0.
 \]
 \item 对 $\Sigma=T^2=\mathbb C/\mathbb Z+i\mathbb Z$，有
 $\dim_{\mathbb C}\operatorname{Aut}(T^2)=1$，所以没有标记点的
 $T^2$ 不稳定。对 $z\in T^2$，
 $\operatorname{Aut}_{\mathbb C}(\Sigma,\{z\})$ 有限（它是什么？），因此 $(\Sigma,z)$ 稳定。
 \item 对亏格 $g(\Sigma)\geq2$ 的黎曼曲面 $\Sigma$，有
 $\dim_{\mathbb C}(\operatorname{Aut}(\Sigma))=0$。否则，从一族连续自同构可得到一个非零全纯向量场 $X$。把 $X$ 看作切丛的截面，并回忆
 $c_1(T\Sigma)=\chi(\Sigma)=2-2g$，计入重数后有
 \[
  \#(\sigma_{T\Sigma}\cap\operatorname{graph}(X))=2-2g,
 \]
 其中 $\sigma_{T\Sigma}$ 是全纯线丛 $T\Sigma\to\Sigma$ 的零截面。因此，当 $g\geq2$ 时，$2-2g<0$。但这导致矛盾，因为 $T\Sigma$ 的任意非零全纯截面 $X$ 与零截面 $\sigma_{T\Sigma}$ 的局部交数总是非负的。
\end{enumerate}
\end{example}

\begin{exercise}\label{ex:automorphisms-three-points-mobius}
证明 $\operatorname{Aut}(\mathbb{CP}^1,z_1,z_2,z_3)=\{id\}$，并写出保持
$\{0\}\subset\mathbb C\subset\mathbb{CP}^1$、
$\{0,\infty\}\subset\mathbb{CP}^1$ 的所有可能分式线性变换。
\end{exercise}

\begin{definition}\label{def:deligne-mumford-stable-curves}
\stfhIndex{I029}{00-symbol-029}{$\overline{\mathcal M}_{g,m}$}
以 $\overline{\mathcal M}_{g,m}$ 表示带 $m$ 个标记点、（算术）亏格为 $g$ 的稳定曲线之同构类集合，即
$g=\operatorname{rank}H^{(1,0)}(\Sigma,\mathbb C)$。以
$\mathcal M_{g,m}$ 表示 $\overline{\mathcal M}_{g,m}$ 中由没有节点的光滑曲线构成的子集。
\end{definition}

空间 $\overline{\mathcal M}_{g,m}$ 由许多层构成，可以借助某类带装饰平面图来分类。现在描述这类图，称为\emph{对偶图}。
\stfhIndex{I029}{00-symbol-029}{$\overline{\mathcal M}_{g,m}$}

\subsection{对偶图}
\label{subsec:dual-graphs-stable-curves}

考虑有限拓扑型的平面图 $T$，即它只有有限条边和有限个顶点。分别以 $V_T,E_T$ 表示顶点集和（无序）边集。注意，边 $e$ 的两个端点 $v_1,v_2$ 可以相同；这对应曲线的一个自交点。

\begin{definition}\label{def:edge-incidence-valence}
\stfhIndex{I048}{00-symbol-048}{$\operatorname{val}(v)$}
\stfhIndex{I298}{jia}{价（valence）}
对 $v\in V_T$，记
\[
 E(v)=\text{与 }v\text{ 关联的边集},
\]
并定义 $V$ 的价为
\[
 \operatorname{val}(v)=\#(E(v)).
\]
当 $e$ 的两个端点都是同一个 $v\in V(T)$ 时，把 $v$ 计数两次。
\end{definition}

\begin{definition}\label{def:decorated-dual-graph}
\stfhIndex{I106}{dui^ou^tu}{对偶图（dual graph）}
\stfhSubIndex{I107}{dui^ou^tu}{对偶图（dual graph）}{qi}{旗（flag）}
\stfhSubIndex{I108}{dui^ou^tu}{对偶图（dual graph）}{qi^jie}{旗阶（flag order）}
\stfhSubIndex{I109}{dui^ou^tu}{对偶图（dual graph）}{kui^ge}{亏格（genus）}
\stfhIndex{I018}{00-symbol-018}{$\operatorname{flag}(v)$}
\stfhIndex{I017}{00-symbol-017}{$\operatorname{flag}(D)$}
定义与三元组 $(T,(g_v),o)$ 相伴的\emph{对偶图} $D$，使下列陈述成立：
\begin{enumerate}[label=(\arabic*)]
 \item $T$ 是满足 $\#(V_T),\#(E_T)<\infty$ 的平面图；
 \item 有函数 $V_T\to\mathbb Z_+$，$v\mapsto g_v$；
 \item 有函数 $o:\{1,\ldots,m\}\to V_T$。
\end{enumerate}
称 $g_v$ 为 $v$ 的\emph{亏格}，称 $o^{-1}(v)$ 为 $v$ 的\emph{旗}。记
$\operatorname{flag}(v):=\#(o^{-1}(v))$，称它为 $v$ 的\emph{旗阶}；称 $m$ 为 $D$ 的旗阶，并记为 $\operatorname{flag}(D)$。
\end{definition}

以 $V_T^{int}\subset V_T$ 表示内部顶点集。提醒读者：我们把没有任何旗附着的顶点看作内部顶点，把旗的末端看作外部顶点。外部边是与一面旗相伴的边。

定义非负整数
\begin{equation}
 g:=\sum_{v\in V_T^{int}}g_v+b_1(T),
 \qquad b_1(T):=\operatorname{rank}H_1(T;\mathbb Z),
 \label{eq:dual-graph-genus}
\end{equation}
称为对偶图 $D$ 的亏格。依定义，
\[
 \operatorname{flag}(D)=\sum_{v\in V_T^{int}}\operatorname{flag}(v).
\]

现在引入两个对偶图之间的同构概念。给出同构定义的最佳方式是采用范畴观点。为此，把 $T$ 中从顶点 $v_1$ 到 $v_2$ 的\emph{边道路}定义为一条从 $v_1$ 出发、到 $v_2$ 结束的定向边线性链。由每个平面图 $T$ 形成一个范畴，其中对象是顶点，态射是边道路。

\begin{exercise}\label{ex:graph-edge-path-category}
证明上述定义使 $T$ 成为一个范畴。
\end{exercise}

以 $\mathcal T$ 表示相应范畴。

\begin{definition}\label{def:dual-graph-isomorphism-combinatorial-type}
如果范畴 $\mathcal T$ 与 $\mathcal T'$ 之间存在一个等价 $\phi$，满足
\[
 g'_{\phi(v)}=g_v\quad\text{对所有 }v\in V_T,
 \qquad \phi\circ o=o',
\]
则称两个对偶图 $D,D'$ 具有相同的组合型，并记为 $D\sim D'$。称函子 $\phi$ 为对偶图之间的同构，称其同构类为图的组合型。
\end{definition}

对每个 $v\in V_T$，定义
\[
 m_v:=\operatorname{val}(v)+\operatorname{flag}(v).
\]

\begin{definition}[稳定性判据]\label{def:stable-dual-graph}
\stfhIndex{I254}{wen^ding^xing^pan^ju}{稳定性判据（stability criterion）}
如果对所有顶点 $v\in V_T$ 都有 $m_v+2g_v\geq3$，则称对偶图
$D=(T,(g_v),o)$ \emph{稳定}。
\end{definition}

现在给出本节的主要定义。

\begin{definition}\label{def:comb-stable-dual-graphs}
\stfhIndex{I008}{00-symbol-008}{$\operatorname{Comb}(g,m)$}
定义 $\operatorname{Comb}(g,m)$ 为亏格为 $g$、旗阶为 $m$ 的稳定对偶图 $D$ 之同构类集合。
\end{definition}

下面的组合结果是一个重要成分，它为固定拓扑型伪全纯映射的不同退化类型数目给出上界。

\begin{theorem}\label{thm:finiteness-stable-dual-graphs}
假设 $T$ 稳定且连通。则有不等式
\[
 \#(V_T)\leq2m+5g-2,
 \qquad \#(E_T)\leq2m+6g-3,
\]
特别地，固定 $g,m$ 时 $D$ 只有有限多个组合型，即
$\#\operatorname{Comb}(g,m)<\infty$。
\end{theorem}

\begin{proof}
先注意
\[
 \#(E_T)=\frac12\sum_{v\in V_T}\operatorname{val}(v).
\]
$T$ 的连通性与欧拉恒等式给出
\stfhIndex{I119}{ou^la^heng^deng^shi}{欧拉恒等式（Euler’s identity）}
\begin{equation}
 \chi(T)=1-b_1(T)=\#V_T-\#E_T
 =\frac12\sum_{v\in V_T}(2-\operatorname{val}(v)),
 \label{eq:euler-identity-dual-graph}
\end{equation}
其中 $\chi(T)$ 是 $T$ 的欧拉示性数。其次，
\begin{equation}
 \sum_vg_v\leq g,
 \qquad b_1(T)\leq g.
 \label{eq:dual-graph-genus-bounds}
\end{equation}
显然，
\[
 \#\{v\in V_T\mid g_v\geq1\}\leq g,
 \qquad
 \#\{v\in V_T\mid\exists e=\{v,v\}\}\leq g,
\]
因为这样的 $e$ 对 $b_1(T)$、从而对 $g$ 有贡献。还有
\[
 \#\{v\in V_T\mid\operatorname{flag}(v)\geq1\}\leq m.
\]
所有其他顶点都满足 $g_v=0$ 且 $\operatorname{flag}(v)=0$，所以稳定性条件蕴含
\[
 3\leq2g_v+m_v
 =2g_v+\operatorname{val}(v)+\operatorname{flag}(v)
 \Longrightarrow\operatorname{val}(v)\geq3.
\]
合起来，已经证明
\begin{equation}
 \#(V_T)\leq m+2g
 +\#\{v\in V_T\mid\operatorname{val}(v)\geq3\}.
 \label{eq:vertex-bound-reduction}
\end{equation}
还需估计 $\#\{v\mid\operatorname{val}(v)\geq3\}$。由\eqref{eq:euler-identity-dual-graph}得到
\[
 \chi(T)\leq\frac12\#\{v\in V_T\mid\operatorname{val}(v)=1\}
 -\frac12\#\{v\in V_T\mid\operatorname{val}(v)\geq3\}.
\]
因此
\begin{equation}
 2\chi(T)+\#\{v\mid\operatorname{val}(v)\geq3\}
 \leq\{v\mid\operatorname{val}(v)=1\}.
 \label{eq:valence-three-versus-one}
\end{equation}
但稳定性条件蕴含
\begin{equation}
 \#\{v\mid\operatorname{val}(v)=1\}
 \leq\#\{v\mid g_v>0\}
 +\#\{v\mid\operatorname{flag}(v)>1\},
 \label{eq:valence-one-stability-bound}
\end{equation}
并且有上界
\begin{equation}
 \#\{v\mid g_v>0\}\leq g,
 \qquad\#\{v\mid\operatorname{flag}(v)>1\}\leq m.
 \label{eq:positive-genus-flag-bounds}
\end{equation}
结合\eqref{eq:dual-graph-genus-bounds}、\eqref{eq:valence-three-versus-one}、\eqref{eq:valence-one-stability-bound}与\eqref{eq:positive-genus-flag-bounds}，得到
\[
 \#\{v\mid\operatorname{val}(v)\geq3\}
 \leq m+g-2\chi(T)
 =m+g-2(1-b_1(T))\leq m+3g-2.
\]
将其代入\eqref{eq:vertex-bound-reduction}，得到
\[
 \#(V_T)\leq2m+5g-2.
\]
再导出
\[
 \#(E_T)=\#(V_T)-(1-b_1(T))\leq2m+6g-3.
\]
有限性断言随即成立，因为旗 $o$ 至多有
\[
 \Pi_m^{\#(V_T)}
 =\binom{\#(V_T)+m-1}{m}
\]
种可能选择。总之，对某个只依赖于 $(g,m)$ 的整数 $N(g,m)$，有上界
\[
 \#(\operatorname{Comb}(g,m))\leq N(g,m).
\]
证明完毕。
\end{proof}

\begin{exercise}\label{ex:estimate-number-combinatorial-types}
用 $g,m$ 估计上述整数 $N(g,m)$。
\end{exercise}

接着在 $\operatorname{Comb}(g,m)$ 上引入偏序。以
$D:=[T,(g_v),o]$ 表示由对偶图 $(T,(g_v),o)$ 代表的同构类。稍微滥用记号，在不会混淆时，偶尔也简写为 $D=(T,(g_v),o)$。

\begin{definition}\label{def:partial-order-dual-graphs}
设 $D,\overline D$ 是 $\operatorname{Comb}(g,m)$ 的两个元素。若
$\overline D=(\overline T,(\overline g_{\overline v}),\overline o)$ 可由
$D:=(T,(g_v),o)$ 经过有限次如下操作得到：把 $T$ 的一个顶点 $v$ 替换为图
$T_v\in\operatorname{Comb}(g_v,m_v)$，并保持图的其余部分不变；则定义
$D>\overline D$。见图 9.2。
\end{definition}

现在按如下方式把每条稳定曲线 $(\Sigma,\mathbf z)$ 与上述带装饰图之一 $D$ 相联系（见图 9.3）：
\begin{itemize}
 \item $T_\Sigma$ 的每个顶点 $v$ 对应 $\Sigma$ 的一个分支 $\{\Sigma_v\}$；
 \item 每条边对应 $\Sigma$ 的一个二重点，每面旗对应一个标记点；
 \stfhIndex{I140}{ji^he^kui^ge}{几何亏格（geometric genus）}
 \item 每个顶点用 $g_v$ 装饰，其中 $g_v$ 是 $\Sigma_v$ 的几何亏格。
\end{itemize}

\begin{figure}[htbp]
 \centering
 \begin{tikzpicture}[x=1cm,y=1cm,>=Stealth,line cap=round,line join=round]
  \begin{scope}[shift={(-3.3,0)}]
    \coordinate (v) at (0,0);
    \draw (-1.35,.8)--(v)--(-1.35,-.8);
    \draw (v)--(1.05,0)--(1.65,.75);
    \draw (1.05,0)--(1.65,-.75);
    \draw (v) .. controls (-.25,.9) and (.65,1.05) .. (.55,.35)
                 .. controls (.45,-.2) and (-.5,-.1) .. (v);
    \draw[dashed] (v) ellipse (.55 and .35);
  \end{scope}
  \draw[->] (-1.0,.55) arc[start angle=150,end angle=30,radius=1.45];
  \node at (0,.95) {替换};
  \begin{scope}[shift={(3.3,0)}]
    \coordinate (a) at (-.7,0); \coordinate (b) at (.7,0);
    \draw (-1.55,.7)--(a)--(-1.55,-.7);
    \draw (b)--(1.55,.7) (b)--(1.55,-.7);
    \draw (a)--(b);
    \draw (a) .. controls (-.65,1.15) and (.1,1.25) .. (.05,.45)
                 .. controls (0,-.05) and (-.75,-.05) .. (a);
    \draw (b) .. controls (.65,1.15) and (1.4,1.25) .. (1.35,.45)
                 .. controls (1.3,-.05) and (.55,-.05) .. (b);
    \draw[dashed] (0,0) ellipse (.95 and .3);
  \end{scope}
\end{tikzpicture}

 \caption{图的偏序。}
 \label{fig:partial-order-graphs}
\end{figure}

\begin{figure}[htbp]
 \centering
 \begin{tikzpicture}[x=1cm,y=1cm,>=Stealth,line cap=round,line join=round]
  \begin{scope}[shift={(-3.2,0)}]
    \coordinate (a) at (0,1); \coordinate (b) at (-.75,.15);
    \coordinate (c) at (.75,.15); \coordinate (d) at (.15,-.8);
    \draw (a)--(b)--(d)--(c)--(a) (b)--(-1.55,-.5) (b)--(-1.3,.85)
          (c)--(1.55,.5) (d)--(.1,-1.65) (d)--(.8,-1.45);
    \draw (d) .. controls (-.75,-.95) and (-.75,-1.6) .. (-.05,-1.45)
              .. controls (.45,-1.3) and (.2,-.9) .. (d);
    \node[above] at (a) {$0$}; \node[left] at (b) {$1$};
    \node[right] at (c) {$1$}; \node[right] at (d) {$2$};
  \end{scope}
  \draw[->] (-.95,0)--(.65,0);
  \begin{scope}[shift={(3.0,0)},scale=.72]
    \draw (-.8,.95) circle (.8); \draw (.8,.9) circle (.65);
    \draw (0,1.95) circle (.58); \draw (0,-.45) ellipse (1.0 and 1.25);
    \draw (-.45,1.45) arc[start angle=200,end angle=340,radius=.55];
    \draw (-.45,1.45) arc[start angle=160,end angle=20,radius=.55];
    \draw (-.55,-.35) .. controls (-.2,.25) and (.45,.1) .. (.55,-.45)
             .. controls (.25,-1.1) and (-.45,-1.05) .. (-.55,-.35);
    \foreach \p in {(-1.35,1.05),(.9,1.0),(-.55,-.8),(.5,-.95),(0,2.35)}
      \fill \p circle (1.7pt);
    \draw[dotted] (-.58,1.8) arc[start angle=180,end angle=360,radius=.58];
    \draw[dotted] (-.75,.05) arc[start angle=180,end angle=360,radius=.75];
  \end{scope}
\end{tikzpicture}

 \caption{图与曲面之间的对应。}
 \label{fig:graphs-surfaces-correspondence}
\end{figure}

于是整数
\[
 g:=\sum_\nu g_v+\operatorname{rank}H_1(T_\Sigma,\mathbb Z)
\]
对应于相伴稳定曲线 $\Sigma$ 的算术亏格（即
$g=\dim H^1(\Sigma,\mathcal O)$，其中 $\mathcal O$ 是 $\Sigma$ 上全纯函数层）。定义
$D_\Sigma=[T_\Sigma,(g_v),o_\Sigma]$，并称 $D_\Sigma$ 或
$(T_\Sigma,(g_v),o_\Sigma)$ 为 $(\Sigma,\mathbf z)$ 的对偶图。在这个对应下，有
\stfhSubIndex{I110}{dui^ou^tu}{对偶图（dual graph）}{z^de}{$(\Sigma,z)$ 的（of $(\Sigma,z)$）}
\[
 \operatorname{val}(v)\longleftrightarrow\operatorname{sing}(v),
 \qquad
 \operatorname{flag}(v)\longleftrightarrow\operatorname{mark}(v),
\]
其中把 $\operatorname{sing}(v)$ 定义为 $\Sigma_v$ 的二重点数，把
$\operatorname{mark}(v)$ 定义为 $\Sigma_v$ 上的标记点数。

反过来，可以把一条稳定曲线 $(\Sigma,\mathbf z)$ 与每个稳定对偶图相联系。

\begin{definition}\label{def:stable-curve-strata-dual-graph}
\stfhIndex{I028}{00-symbol-028}{$\mathcal M^D_{g,m}$}
\stfhIndex{I029}{00-symbol-029}{$\overline{\mathcal M}_{g,m}$}
以 $\mathcal M_{g,m}^D$ 表示对偶图为 $D$ 的稳定曲线之同构类集合，并作并集
\[
 \overline{\mathcal M}_{g,m}
 :=\coprod_{D\in\operatorname{Comb}(g,m)}\mathcal M_{g,m}^D.
\]
\end{definition}

\begin{exercise}\label{ex:stable-curves-second-countable-locally-compact}
证明关于稳定曲线等价类上自然诱导的拓扑，$\mathcal M_{g,m}$ 是第二可数且局部紧的。
\end{exercise}

于是可以把 $\mathcal M_{g,m}$ 表示为可数并
\[
 \mathcal M_{g,m}=\bigcup_{\ell=1}^\infty\mathcal U_\ell,
 \qquad \mathcal U_\ell\subset\overline{\mathcal U}_\ell
 \subset\mathcal U_{\ell+1},
\]
且每个 $\overline{\mathcal U}_\ell$ 都紧。称由真映射
$\ell\longmapsto[\Sigma_\ell,j_\ell]$（$\ell\in\mathbb N$）给出的
$\mathcal M_{g,m}$ 中任意序列为\emph{退化序列}。

下面命题的证明留作习题。

\begin{proposition}\label{prop:two-kinds-stable-curve-degeneration}
假设 $[\Sigma_\ell,j_\ell]\in\mathcal M_{g,m}$ 是退化序列。则有两类退化：
\begin{enumerate}[label=(\roman*)]
 \item 两个特殊点的碰撞；
 \item 共形结构的退化。
\end{enumerate}
\end{proposition}

\begin{exercise}\label{ex:two-kinds-degeneration-proof}
证明这个命题（细节也可见 (DM69)）。（提示：证明其逆否命题。先用庞加莱度量，即唯一双曲度量，表示相伴带孔黎曼曲面的一个共形结构；再证明，若共形结构不发生退化，则单射半径一致地远离零。注意稳定性蕴含相伴庞加莱度量是双曲的。）
\end{exercise}

更几何地说，就 $\Sigma\setminus\{z_1,\ldots,z_m\}$ 上的庞加莱度量而言，带标记黎曼曲面的退化对应于夹断庞加莱度量的一条短测地线。因此，除非所得曲线所围圆盘的内部至少包含两个标记点，否则不能夹断一条同调平凡环路。这是因为，不然的话，由短测地线围成的这种区域不能带有常曲率 $-1$ 的度量。

现在希望在带标记稳定曲线集合 $\overline{\mathcal M}_{g,m}$ 上赋予一个反映带标记黎曼曲面这些退化的拓扑。更准确地说，将给出满足下列条件的拓扑：
\begin{enumerate}[label=(\arabic*)]
 \item 它在 $\mathcal M_{g,m}\subset\overline{\mathcal M}_{g,m}$ 上的子空间拓扑等价于原商拓扑
 $\mathcal M_{g,m}=\widetilde{\mathcal M}_{g,m}/\sim$；
 \item 每个退化序列都导致在定义 9.3.16 的偏序中下降到一个对偶图。
\end{enumerate}

\begin{figure}[htbp]
 \centering
 \begin{tikzpicture}[x=1cm,y=1cm,line cap=round,line join=round]
  \coordinate (N) at (0,1.35); \coordinate (E) at (2.05,0);
  \coordinate (S) at (0,-1.35); \coordinate (W) at (-2.05,0);
  \foreach \p/\lab/\pos in {N/{(a)}/above,E/{(a)}/right,S/{(a)}/below,W/{(a)}/left}{
    \draw (\p) circle (.65); \node[\pos] at ($(\p)+(0,0)$) {\lab};
  }
  \draw ($(W)+(.35,.43)$) .. controls (-1.15,.65) and (-.68,.82) .. ($(N)+(-.42,-.35)$)
        -- ($(N)+(0,-.75)$)
        .. controls (-.78,.28) and (-1.25,.25) .. ($(W)+(.35,-.43)$);
  \draw ($(N)+(.42,-.35)$) .. controls (.68,.82) and (1.15,.65) .. ($(E)+(-.35,.43)$)
        -- ($(E)+(-.75,0)$)
        .. controls (.78,.28) and (.78,-.28) .. ($(E)+(-.35,-.43)$);
  \draw ($(E)+(-.35,-.43)$) .. controls (1.15,-.65) and (.68,-.82) .. ($(S)+(.42,.35)$)
        -- ($(S)+(0,.75)$)
        .. controls (.78,-.28) and (1.25,-.25) .. ($(E)+(-.35,.43)$);
  \draw ($(S)+(-.42,.35)$) .. controls (-.68,-.82) and (-1.15,-.65) .. ($(W)+(.35,-.43)$)
        -- ($(W)+(.75,0)$)
        .. controls (-.78,-.28) and (-.78,.28) .. ($(W)+(.35,.43)$);
\end{tikzpicture}

 \caption{归纳过程。}
 \label{fig:inductive-procedure-orbifold-charts}
\end{figure}

$\overline{\mathcal M}_{g,m}$ 上这种拓扑的基本要求如下：
\begin{enumerate}[label=(\arabic*)]
 \item 它是紧且豪斯多夫的；
 \item 若以 $\overline{\mathcal M}_{g,m}^D$ 表示
 $\mathcal M_{g,m}^D$ 在 $\overline{\mathcal M}_{g,m}$ 中的闭包，则
 \[
  \overline{\mathcal M}_{g,m}^D
  =\coprod_{D'\leq D}\mathcal M_{g,m}^{D'};
 \]
 \item 它具有坐标卡，使坐标邻域同胚于 orbifold $V/\Gamma$，其中
 $V\subset\mathbb R^N$ 是子空间，$\Gamma$ 是
 $GL(\mathbb R^N)$ 的有限子群（见图 9.4）；
 \item 存在普适丛
 $\mathcal U_{g,m}\to\overline{\mathcal M}_{g,m}$，其在
 $[\Sigma,\mathbf z]\in\overline{\mathcal M}_{g,m}$ 上的纤维同构于稳定曲线自身 $(\Sigma,\mathbf z)$。
\end{enumerate}

事实上，有构成局部平凡纤维化（在 orbifold 意义下）的自然遗忘映射
\begin{equation}
\begin{split}
 \operatorname{forget}_i:\mathcal M_{g,m+1}&\longrightarrow\mathcal M_{g,m},\\
 [j,(z_1,\ldots,z_{m+1})]&\longmapsto
 [j,(z_1,\ldots,\widehat z_i,\ldots,z_{m+1})].
\end{split}
 \label{eq:forgetful-map-smooth-marked-curves}
\end{equation}
它的纤维
$\operatorname{forget}_i^{-1}([j,(z_1,\ldots,z_m)])$ 可等同于
\[
 \bigl\{(j',(z_1',\ldots,z,\ldots,z_m'))\mathrel{\big|}
 z\in\Sigma\setminus\{z_1',\ldots,z_m'\}\bigr\}/
 \operatorname{Aut}(\Sigma,\mathbf z'),
\]
它同构于
\[
 \{z\in\Sigma\setminus\{z_1,\ldots,z_m\}\}/
 \operatorname{Aut}(j,(z_1,\ldots,z_m)).
\]

这个遗忘映射连续延拓到普适纤维化
$\mathcal U_{g,m}\to\overline{\mathcal M}_{g,m}$；后者本身由遗忘映射
\begin{equation}
 \operatorname{forget}_{m+1}:
 \overline{\mathcal M}_{g,m+1}\longrightarrow
 \overline{\mathcal M}_{g,m}
 \label{eq:forgetful-map-stable-curves}
\end{equation}
表示。

\begin{remark}\label{rem:stabilization-after-forgetting-mark}
与\eqref{eq:forgetful-map-smooth-marked-curves}不同，遗忘映射
\eqref{eq:forgetful-map-stable-curves}的定义需要一些解释。遗忘最后一个，即第 $(m+1)$ 个标记点后，包含该标记点的不可约分支可能变得不稳定。这只可能发生在该分支亏格为零且恰有三个特殊点时。有两种情形。一种情形是三个特殊点由两个节点和一个标记点组成；另一种情形是它们由一个节点和一个标记点组成。

在第一种情形，只需压缩不稳定分支。容易看出，所有其他分支都不受这个过程影响，因此所得预稳定曲线是稳定的。在第二种情形，压缩该分支，并在相邻分支上被压缩分支扎根的位置留下一个剩余标记点。所得预稳定曲线同样稳定。此外，标记点总数不变，亏格也保持不变，因为压缩的是亏格 $0$ 分支。这就是遗忘映射\eqref{eq:forgetful-map-stable-curves}的含义。
\end{remark}

还有下面的一般引理，其证明留作习题。

\begin{lemma}\label{lem:forgetful-map-compositions}
在 $\overline{\mathcal M}_{g,m+n+n'}$ 中，
\[
 \operatorname{forget}_{m+n;m}\circ
 \operatorname{forget}_{m+n+n';m+n}
 =\operatorname{forget}_{m+n+n';m}
 =\operatorname{forget}_{m+n';m}\circ
 \operatorname{forget}_{m+n+n';m+n'}.
\]
\end{lemma}

\begin{exercise}\label{ex:forgetful-map-compositions}
证明这个引理。
\end{exercise}
